\input{preamble}

\begin{document}

\title{Algebra}
\maketitle

\phantomsection
\label{section-phantom}
\hfill
\href{http://github.com/danimalabares/stack}{github.com/danimalabares/stack}

\tableofcontents

\section{Rings}
\label{section-rings}

\noindent
So elementary that this is not in Stacks Project.

\begin{theorem}
\label{theorem-ideals-quotient-correspondence}
Let $R$ be a ring and $I \subset R$ an ideal.
There is a correspondence between ideals of $R$ containing $I$ 
and ideals of $R/I$.
\end{theorem}

\begin{proof}
We show that there are two maps between the set of ideals of $R$ 
containing $I$ and the ideals of $R/I$ whose compositions
are the identities in the corresponding set. One map is the projection
to the quotient, and the other map is taking the elements
whose equivalence class lies in a given ideal.

Suppose that $\tilde{J}$ is an ideal of $R/I$.
Define $J=\{j \in R : j+I \in \tilde{J}\}$, 
the set of elements that project to $\tilde{J}$.
Then $I \subset J$, because $i + I = I$ is the zero of $R/I$ 
which is contained in any ideal of $R/I$.
Also notice that $J$ is an ideal of $R$ since for any $r \in R$,
$rj+I \in \tilde{J}$ by $\tilde{J}$ being an ideal.
Notice that projecting $J$ back to $R/I$ gives $J$
by definition.

Conversely, we project to the quotient.
Let $J$ be any ideal of $R$ containing $I$ 
Let $\tilde{J}$ be the projection to $R/I$.
Then $\tilde{J}$ is an ideal of $R/I$ by $J$ being an ideal of $R$.
Consider the set $J'$ of all elements whose projection lies in $\tilde{J}$.
Then $J'=J$ by definition.
\end{proof}

\noindent
This is the essential tool in proving that $R/I$ is a field
when $I$ is a maximal ideal

\begin{lemma}
\label{lemma-quotient-ring-maximal-ideal-is-field}
Let $R$ be a ring and $I$ a maximal ideal of $R$.
Then $R/I$ is a field.
\end{lemma}

\begin{proof}
It's easy to prove that a ring is a field if and only if its
only ideals are $\{0\}$ and $R$. By Theorem \ref{theorem-correspondence-ideals},
the ideals of $R/I$ are in correspondence with ideals of $R$ containing $I$,
but only $R$ is such.
\end{proof}

\section{Fields}
\label{section-fields}

\noindent
Just for the record, any simple field extension, i.e. the smallest
field in $F$ containing a single element, is isomorphic
to either que ``rational function field $k(t)/k$''
or to one of the field extensions $k[t]/(P)$ 
with $P \in k[t]$ irreducible.
(Stacks Project tag \href{https://stacks.math.columbia.edu/tag/09G1}{09G1}.)

\section{Modules}
\label{section-modules}

It looks like most books (at least
\cite{eis} and \cite{Samuel-Zariski-Vol1})
define a module to be an abelian group
along with a certain operation with
elements of the ring.

Which secretly just says that

\begin{definition}
\label{definition-module}
Let $R$ be a ring.
An {\it (left) $R$-module} $M$ is an abelian group
such that there exists a (left) ring morphism
$$
R \to \text{End}(M).
$$
\end{definition}

\noindent
Indeed, the three usual requirements
correspond to:

\begin{enumerate}
\item The endomorphism corresponding to $a \in R$ 
respect the group structure of the group $M$:
$$
a(x+y)=ax+ay,
$$

\item The representation map is a map of rings
$$
(a+b)x=ax+by,\qquad (ab)x=a(bx).
$$
\end{enumerate}

\section{Tensor product}
\label{section-tensor-product}

\begin{definition}
\label{definition-tensor-product}

\end{definition}

\begin{theorem}[Universal property of tensor product]
\label{theorem-universal-property-of-tensor-product}
Let $M_1,M_2,M$ be modules over a ring $R$.
For any bilinear map $B:M_1\times M_2 \to M$ there exists
a unique homomorphism $b:M_1 \otimes M_2 \to M$
making the following diagram commute
$$
\xymatrix{
M_1\times M_2 \ar[r]^\pi\ar[dr]_B &  M_1 \otimes M_2 \ar[d]^b\\
& M
}
$$
where $\pi(m_1,m_2)=m_1 \otimes m_2$.
\end{theorem}

\begin{remark}
\label{remark-whitneys-definition}
For $R=k$ the above property shows that
$\text{Bil}(M_1\times M_2,k)=(M_1 \otimes M_2)^*$.
If we assume that the vector spaces are finite dimensional,
we can take duals to obtain
$$
M_1 \otimes M_2 = (M_1 \otimes M_2)^ {* *}=\text{Bil}(M_1 \times M_2,k)^*,
$$
which is how Whitney first defined tensor products.
\end{remark}

\noindent
Here's Misha's trick on how to show that the universal property
of tensor product implies its uniqueness.
Consider the category $\mathcal{C}$ of pairs of an $R$-module $M$,
and a bilinear map $M_1 \times M_2 \to M$.
Then we prove that the tensor product $M_1 \otimes M_2$
is an initial object in such category,
and it is easy to show that initial objects are unique.

\begin{remark}
\label{remark-universal-property}
Misha's alternative way of stating the universal property of tensor product is
that there is a bijective correspondence
$$
\Hom_R(M_1 \otimes_R M_2, M)
\xymatrix{\ar@{<~>}[r]&}
\text{Bil}(M_1\times M_2,M).
$$
The key observation is that the LHS are {\it linear} maps from a module to $M$
while the RHS are {\it bilinear} maps.
\end{remark}

\begin{exercise}[Currying]
\label{exercise-currying}
$$
\Hom(M_1\otimes_R M_2, M)
=
\Hom(M_1,\Hom(M_2,M))
$$
\end{exercise}

\begin{proof}
This is implied by the universal property
since the RHS is just $\text{Bil}(M_1\times M_2,M)$.
\end{proof}

\begin{exercise}
\label{exercise-left-hom-sends-right-exact-to-left-exact}
Let 
$$
\xymatrix{
M_1\ar[r]^{f_1}
&M_2\ar[r]^{f_2}
&M_3\ar[r]
&0
}
$$
be exact.
Show that
$$
\xymatrix{
0\ar[r]
&\Hom(M_3,N)\ar[r]^{-\circ f_2}
&\Hom(M_2,N)\ar[r]^{- \circ f_1}
&\Hom(M_1,N)
}
$$
is exact.
\end{exercise}

\begin{proof}
Injectivity of the first arrow is easy.
For the middle term pick $f \in \Hom_R(M_2,N)$
in the kernel of $- \circ f_1$.
To show that it is in the image of $-\circ f_2$
we look for $g \in \Hom_R(M_3,N)$
such that $f=g\circ f_2$.

Since $M_1 \to M_2 \to M_3 \to 0$
is exact, we have an isomorphism
$\bar{f}_2:M_2/\text{img}f_1 \xrightarrow{\simeq} M_3$.
Since $f \circ f_1=0$, we see that $f$
induces a map $\bar{f}$ on the quotient since
Indeed, if $m_2-m_2' \in \text{img}f_1$,
we have $0=f(m_2-m_2')=f(m_2)-f(m_2')$.

Then for any $m_3 \in M_3$
we have a unique class $m_2 + M_1 \in M_2/\text{img}f_1$.
Thus we can define $g = \bar{f} \circ \bar{f}_2^{-1}$.

Note that this choice satisfies $f = g \circ f_2$.
Indeed,
$$
g\circ f_2(m_2)=\bar{f} \circ \bar{f}^{-1}_2(m_3)
= \bar{f} \circ (m_2+M)
=f(m_2).
$$
\end{proof}

\begin{exercise}
\label{exercise-counterexample-right-hom-is-not-right-exact}
Let $R = \mathbb{Z}$.
Find an example of an injective homomorphism
$M_1 \to M_2$ such that $\Hom(M_2,N) \to \Hom(M_1,N)$
is not surjective.
\end{exercise}

\begin{proof}
Consider
$$
\xymatrix{
0\ar[r]
&2\mathbb{Z}\ar[r]^i\ar[d]_\varphi
&\mathbb{Z}\ar@{-->}[dl]^{\not\exists}\\
&\mathbb{Z}/2\mathbb{Z}
}
$$
where
\begin{align*}
\varphi: 2\mathbb{Z} &\longrightarrow \mathbb{Z}/2\mathbb{Z} \\
2n &\longmapsto [n].
\end{align*}

\noindent
Then there is no extension $\bar{\varphi}:\mathbb{Z} \to \mathbb{Z}/2\mathbb{Z}$
making $\varphi=\bar{\varphi} \circ i$.
Indeed, note that $\varphi(2)=[1]$,
while $\bar{\varphi}(2\cdot 1)=2\bar{\varphi}(1)=[0]$.
\end{proof}

\noindent
The idea is that any element in the image must be divisible
by any integer. This is why $\mathbb{Q}$ is an injective module
according to the following definition.
(A precise proof of why $\mathbb{Q}$ is injective
can be found at
\href{https://stacks.math.columbia.edu/tag/01D6}{01D6}.)

\begin{definition}
\label{definition-injective-module}
An $R$-module is {\it injective} if the functor $\Hom(-,M)$
is exact.
That is, if the dotted arrow in the following diagram always exists:
$$
\xymatrix{
0\ar[r]
&M_1\ar@{^{(}->}[r]\ar[d]
&M_2\ar@{-->}[dl]^{\exists}\\
&N
}
$$
\end{definition}

\begin{exercise}
\label{exercise-left-hom-is-left-exact}
Let
$$
\xymatrix{
0\ar[r]
&M_1\ar[r]^{f_1}
&M_2\ar[r]^{f_2}
&M_3\ar[r]
&0
}
$$
be an exact sequence of $R$-modules.
Prove that
$$
\xymatrix{
0\ar[r]
&\Hom_R(N,M_1)\ar[r]^{f_1 \circ -}
&\Hom_R(N,M_2)\ar[r]^{f_2 \circ -}
&\Hom_R(N,M_3)
}
$$
is exact
\end{exercise}

\begin{proof}
Injectivity is simple.
For the middle term let $f:N \to M_2$ be in the kernel
of $f_2 \circ -$,
that is, $f(N)$ is killed by $f_2$.

We want to define a map $g:N \to M_1$
such that $f_1 \circ g = f$.
But observe that the image of $f$ already lies in $M_1 = \Ker f_2$.
Thus $f$ is already a map $f:N \to M_1$,
save a more precise notation.
\end{proof}

\begin{exercise}
\label{exercise-counterexample-left-hom-not-right-exact}
Let $R=\mathbb{Z}$. Find an epimorphism
(surjective homomorphism)
of $R$-modules $M_2 \to M_3$ such that the corresponding map
$$
\Hom_R(N,M_2) \to \Hom_R(N,M_3)
$$
is not surjective for some $R$-module $N$.
\end{exercise}

\begin{proof}
Now consider
$$
\xymatrix{
&\mathbb{Z} \ar[d]^\pi\ar@{-->}[dl]_{\not\exists}\\
2\mathbb{Z}\ar[r]_-\varphi& \mathbb{Z}/2\mathbb{Z}\ar[r]& 0
}
$$
for the same $\varphi$ as in Exercise
\ref{exercise-counterexample-right-hom-is-not-right-exact}.
Then $\pi(1)=[1]$, while any extension mapping $1 \to 2n$
gives $\varphi(2n)=2\varphi(n)=[0]$.
\end{proof}

\begin{definition}
\label{definition-projective-module}
An $R$-module $N$ is {\it projective} if the functor
$\Hom_R(N,-)$ is exact. That is, if the dotted
arrow in the following diagram always exists:
$$
\xymatrix{
&N \ar[d]\ar@{-->}[dl]_{\exists}\\
M_2\ar@{->>}[r]&M_3\ar[r]& 0
}
$$

\end{definition}

\begin{exercise}
\label{exercise-finitely-generated-projective-iff-factor-of-free}
A finitely generated $R$-module $N$ is projective
if and only if
it is a direct summand of a free $R$-module $R^n$,
i.e. there is a direct sum decomposition
$R^n=N \oplus N'$.
\end{exercise}

\begin{proof}
Consider $R^n \simeq \bigoplus_{i=1}^n a_iR$
where $N$ is generated by $a_1,…,a_n$.
Apply projectivity to obtain
$$
\xymatrix{
&N \ar[d]^{\text{id}}\ar@{-->}[dl]_{\exists \psi}\\
R^n\ar[r]_-\varphi& N \ar[r]& 0
}
$$
where
\begin{align*}
\varphi: R^n &\longrightarrow N \\
(r_1a_1,…,r_na_n) &\longmapsto \sum r_ia_i.
\end{align*}

\noindent
This is a section of $\varphi$ in the exact sequence
$$
\xymatrix{
0\ar[r]
&\Ker \varphi\ar[r]
&R^n\ar[r]^{\varphi}
&N\ar[r]\ar@/_-.5pc/[l]^\psi
&0
}
$$
which gives decomposition
$$
R^n \simeq \Ker\varphi \oplus \text{Im}\psi.
$$
More explicitly, the intersection of those modules is clearly trivial
while any element of $R^n$ may be written as
$$
x=\underbrace{(x-\psi \varphi(x))}_{\in \Ker \varphi}
+\underbrace{\psi\varphi(x)}_{\in \text{Im}\psi}.
$$
For the converse we start from the diagram given in the definition
of projective module, add $N'$ to every term, arriving 
at $R^n\simeq N\oplus N'$. We observe that free modules
are projective by constructing an extension using a basis of $R^n$,
and then back to $N$ by taking the action of the extension restricted to $N$.
\end{proof}

\medskip\noindent
In fact, currying (Exercise \ref{exercise-currying}) 
is the essential ingredient
for the proof that tensor product is left exact.
Such result follows from two lemmas.

\begin{lemma}
\label{lemma-tensor-product-functor-is-left-exact}
Let 
$$
\xymatrix{
0\ar[r]
&M_1\ar[r]
&M_2\ar[r]
&M_3\ar[r]
&0
}
$$
be an exact sequence of $R$-modules.
Then for any $R$-modules $N$ and $N'$,
the sequence
$$
\xymatrix{
0\ar[r]
&\Hom_R(M_3 \otimes N',N)\ar[r]
&\Hom_R(M_2 \otimes N',N)\ar[r]
&\Hom_R(M_1 \otimes N',N)
}
$$
is exact.
\end{lemma}

\begin{proof}
Apply $\Hom_R(-,N)$,
then $\Hom_R(N',-)$
and then currying.
\end{proof}

\noindent
There is an (almost) converse 
to Exercise \ref{exercise-left-hom-sends-right-exact-to-left-exact}:

\begin{lemma}
\label{lemma-converse-hom-is-left-exact}
Let
$$
\xymatrix{
E:&
M_1\ar[r]^\mu
&M_2\ar[r]^\rho
&M_3\ar[r]
&0
}
$$
be a complex,
i.e. $\rho\mu=0$,
such that
$$
\xymatrix{
0\ar[r]
&\Hom_R(M_3,N)\ar[r]^{\rho_N}
&\Hom_R(M_2,N)\ar[r]^{\mu_N}
&\Hom_R(M_1,N)
}
$$
is exact for any $R$-module $N$. Then $E$ is also exact.
\end{lemma}

\noindent
And this is what implies right-exactness of tensor product:

\begin{theorem}
\label{theorem-tensor-product-is-right-exact}
Let
$$
\xymatrix{
M_1\ar[r]
&M_2\ar[r]
&M_3\ar[r]
&0
}
$$
be exact. Then
$$
\xymatrix{
M_1 \otimes_R M\ar[r]
&M_2 \otimes_R M\ar[r]
&M_3 \otimes_R M\ar[r]
&0
}
$$
is exact.
\end{theorem}

\begin{proof}
This is just Lemma \ref{lemma-tensor-product-functor-is-left-exact}
and Lemma \ref{lemma-converse-hom-is-left-exact}!!!
\end{proof}

\begin{exercise}
\label{exercise-not-flat}
Let $R=\mathbb{C}[t]$.
Find an exact sequence of $R$-modules
$$
\xymatrix{
0\ar[r]
&M_1\ar[r]
&M_2\ar[r]
&M_3\ar[r]
&0
}
$$
and an $R$-module $N$
such that the corresponding map
$$
M_1 \otimes_R N \to M_2 \otimes_R N
$$
is not injective.
\end{exercise}

\begin{proof}
Suppose the given exact sequence is just the ideal
exact sequence of a point in $\mathbb{A}^1$.
Then tensoring by $M$ means thickening the point.
Not having injectivity after tensoring
means that tensoring the ideal by $M$
does not give the ideal of the resulting variety.
\end{proof}

And from this it follows as corollary

\begin{proposition}
\label{proposition-tensor-ideals}
Let $I \subset R$ be an ideal of a ring.
Then $M \otimes_R (R/I)=M/IM$.
\end{proposition}

\begin{proof}
Take the exact sequence $0 \to I \to R \to R/I$ and tensor
by $M$.
\end{proof}

\begin{exercise}
\label{exercise-tensor-product-of-quotients-by-maximal-ideals-is-zero}
Let $I,I'$ be distinct maximal ideals in a ring $R$.
Then $R/I \otimes_R R/I'=0$.
\end{exercise}

\begin{proof}
They key observation is that since $I$ and $I'$
are maximal we must have $I+I'=M$.
This means that $1=f+f'$ for some $f \in I$ and $f' \in I'$.
Then
$1(a+I)\otimes(b+I')
=(fa+I) \otimes (b + I')
+(a+I)\otimes(fb+I')=0$.
\end{proof}



\section{Algebras}
\label{section-algebras}

\noindent
There appears to be no definition of ``algebra''
in Stacks Project.
But there is one in \cite{Eisenbud}:

\begin{definition}
\label{definition-algebra}
If $R$ is a commutative ring,
then a {\it commutative algebra} over $R$ 
is a commutative ring $S$ together
with a ring morphism $R \to S$.
\end{definition}

\noindent
But actually I was expecting that 
$S$ would be defined as 
a ring that is also an $R$-module.
So let us note that a morphism of rings
$R \to S$ gives a representation
$R \to \text{End}(S)$ via left multiplication.
But the other way around,
given an endomorphism associated
to some element in $R$,
how do we assign an element of $S$
so as to produce a map $R \to S$?

\section{Finitely-generated and finitely presented algebras}
\label{section-finitely-presented-generated}

\noindent
See Stacks Project tag \href{https://stacks.math.columbia.edu/tag/00F2}{00F2}.

Upshot: a finitely-generated $R$-algebra $S$ is such that
$S \simeq R[x_1,\ldots,x_n]/I$ for some ideal $I \subset R$.
Another way of saying this is that there is a surjective morphism
$R[x_1,…,x_n]\to S$.
Finitely-presentedness is when $I = (s_1,\ldots,s_k)$.

As an aside, recall that finitely presented $k$-algebras
with no nilpotents are in correspondence with affine algebraic varieties.

\begin{definition}
\label{definition-finite-type}
Let $R \to S$ be a ring map.
\begin{enumerate}
\item We say $R \to S$ is of {\it finite type}, or that {\it $S$ is a finite
type $R$-algebra} if there exist an $n \in \mathbf{N}$ and a surjection
of $R$-algebras $R[x_1, \ldots, x_n] \to S$.
\item We say $R \to S$ is of {\it finite presentation} if there
exist integers $n, m \in \mathbf{N}$ and polynomials
$f_1, \ldots, f_m \in R[x_1, \ldots, x_n]$
and an isomorphism of $R$-algebras
$R[x_1, \ldots, x_n]/(f_1, \ldots, f_m) \cong S$.
\end{enumerate}
\end{definition}

IA: for modules, the definition of finitely generated $R$-module
is that there exists a surjective morphism $R^n \xymatrix{\ar@{->>}[r]&} S$.
So: f.g. algebra means elements are polynomials 
on the generators with coefficients in the ring,
and f.g. module means elements are linear combinations
on the generators with coefficients in the ring.

\section{Noetherian and Artinian rings}
\label{section-noetherian-and-artinian-rings}

\begin{definition}
\label{definition-noetherian-and-artinian-rings}
A ring is called {\it Noetherian} 
if any increasing chain of ideals $I_1 \subset I_2 \subset…$
stabilizes. It is called {\it Artinian} 
if any descending chain of ideals $I_1 \supset I_2\supset…$
stabilizes.
\end{definition}

\begin{exercise}
\label{exercise-noetherian-iff-ideals-are-finitely-generated}
A ring is Noetherian if and only if all its ideals are finitely generated.
\end{exercise}

\begin{proof}
First we prove that if not every ideal is finitely generated
then $R$ cannot be Noetherian.
Indeed, if $I$ is not finitely generated
we construct an ascending chain by of ideals by starting with
any generator of $I$ and adding other generators one by one.

For the converse suppose that all ideals of $R$ are finitely generated
and pick an ascending chain $I_1\subset I_2\subset…$.
Consider $I=\bigcup_i I_i$,
which is an ideal since for any $a \in I$
there is $i$ such that $a \in I_i$
so that $ra \in I_i$ for any $r \in R$.
However, since it is finitely generated,
all the generators have to be in some $I_j$,
and thus the chain stabilizes.
\end{proof}

\begin{theorem}[Hilbert basis theorem]
\label{theorem-hilbert-basis}
If $R$ is Noetherian then so is $R[x]$.
\end{theorem}

\begin{exercise}
\label{exercise-one-prime-ideal-is-artinian-or-not}

\end{exercise}

\begin{proof}
$k[x_1,x_2,…]/(x_1^2,x_2^2,…)$,
$I_n=(x_n,x_{n+1},…)$
then $I_1 \supset I_2 \supset …$.
\end{proof}

\begin{exercise}
\label{exercise-hilbert-basis-for-formal-power-series}
If $A$ is Noetherian then so is $A[[t]]$.
\end{exercise}

\begin{proof}
Pass from an ideal in $A[[t]]$ to $A$ by taking
the coefficients of the monomials of minimal degree
of each element in the ideal.
Then take a finite set of generators.
Then go back to $A[[t]]$ by taking the polynomials
where the generators where taken.
Then…
\end{proof}

\section{Localization}
\label{section-localization}

\noindent
The general recipe for localization is that
$\mathcal{O}_U=\mathcal{O}_X[f^{-1}]$ where
$U=X\setminus V(f)$.
That is, localizing gives the sheaf of the variety without the zero set.

\begin{definition}
\label{definition-localization}
Let $R$ be a ring and $S \subset R$ 
a multiplicative subset of $R$ (containing $1$).
The {\it localization} of $R$ at $S$ 
is the ring $RS^{-1}$ 
(where $S^{-1}$ is just a set with
the same elements as $S$ but with a $-1$ exponent)
with the operations of $(rs^{-1})(r_1s_1^{-1})=rr_1s^{-1}s_1^{-1}$,
$rs^{-1}+r_1s_1^{-1}=(rs_1+r_1s)(s^{-1}s^{-1}_1)$
and the identification given by
$rs^{-1}=r_1s_1^{-1} \iff rs_1=r_1s$.
(Check this definition if $R$ has zero divisors!)
\end{definition}

\noindent
In a simple case, consider a ring $R$ and $f \in R$.
The {\it localization of $R$ at $f$}
is $R[t]/(ft-1):=R[f^{-1}]=R(f)$. 
The quotient basically says that $t$ is the inverse of $f$.

\begin{example}
\label{example-localization}
\begin{itemize}
\item Dyadic numbers: $\mathbb{Z}[2^{-1}]$,
which are fractions of integers with denominators which
are powers of $2$.

\item Laurent polynomials: $\mathbb{C}[t,t^{-1}]$.
\end{itemize}
\end{example}

\begin{lemma}
\label{lemma-localization-nilpotent}
If $R[f^{-1}]=0$ then $f$ is nilpotent.
\end{lemma}

\begin{proof}
If $R[f^{-1}]=0$ then $1 \in (ft-1)$,
which means that $1=(1-ft)p(t)$ for some $p(t) \in k[t]$.
Then…
\end{proof}

\noindent
The localization a ring is not always a local ring.
But if $S$ is a prime ideal, then it is.


\begin{exercise}
\label{exercise-localization-of-noetherian-is-noetherian}
Prove that localization of a Noetherian ring is Noetherian.
\end{exercise}

\begin{proof}
Let $R$ be a Noetherian ring and $S \subset R$.
Let $I_1\subset I_2 \subset…$
be a chain of ideals in the localization $R[S^{-1}]$.
Let $\pi:R \to R[S^{-1}]$ the canonical projection $r \mapsto r$.

Let $I$ be any of the $I_i$.
Consider $\pi^{-1}(I)$. This gives a chain of ideals in $R$,
which must stabilize.
To go back to the localization simply recall that
$\pi^{-1}(I)[S^{-1}]=I$.
\end{proof}





\section{Finite and integral ring extensions}
\label{section-finite-ring-extensions}

\noindent
The upshot is: let $\varphi:R \to S$ be a ring map.
An element $s \in S$ is integral over $R$ if there
is a monic polynomial with coefficients in $R$
that gives zero when you put $s$ instead of the variable.

But it might be convenient to think about
finite (or, for humans, finitely-generated) $R$-modules:
$x$ is integral over $R$ if and only if $R[x]$
is a finitely-generated  $R$-module (i.e.,
not only finitely generated as an algebra
but also as a module).
See \href{
https://youtu.be/5fCHj80BQHw?si=o7tMC4tRsbBnwmJi}{Zvi Rosen's video}.
The forward implication is: suppose $x$ is integral,
then $x^n+r_1x^{n-1}+\ldots+r_n=0$, so
$x^n=-r_1x^{n-1}+\ldots+r_n$,
which ``says'' $R[x] \cong \bigoplus_{k=0}^{n-1}Rx^k$.

So maybe the details of that are not entirely trivial (Zvi uses some theorem,
and Stacks Project does lots of things).
But the point is: I think that after applying Spec,
the fibers of an integral extension are {\bf finite}.
Which of course has all to do with the finite-moduleness
I was trying to explain.




\section{Gröbner theory}
\label{section-grobner}

\noindent
So far here's some rather unformatted notes from
Lin's talk at Commutative Algebra school in Recife 2025.

Would SAGBI bases be related to this?

square free monomial ideals preserve regularity upon Grobner deformation
i.e. initial ideal!

Constantini-Fouli-Goel-Lin-Lindo-Liske-Mostafazadehfard, 2025.

https://sites.psu.edu/kul20/

How is Lin using the Reese algebra here?
Is it ture that the initial ideal of the Gröbner deformation
coincides with the initial ideal ``of the Resse algebra''?



\bibliography{my}
\bibliographystyle{amsalpha}

\end{document}
